\section{Results}
\label{sec:results}

\subsection{Planned Analysis}
\label{sec:planned-analysis}
% This subsection is a pre-registered analysis plan, fixed at design time
% before any model output exists, precisely so that analysis choices
% aren't made in response to whatever pattern the confirmatory run
% happens to show. The actual results (below) are blocked on the
% confirmatory run and are drafted separately once that data exists.

For RQ1--RQ2, \texttt{fault\_rating} is modeled with a mixed-effects
specification with agent gender configuration, severity, and their
interaction with norm family as fixed effects, and \emph{scenario} (the
36 base narratives, prior to gender/severity rendering) as a random
effect -- this is the reason scenario was elevated to a real manipulated
dimension with within-family replication (Section~\ref{sec:dataset})
rather than one hand-written template per family: a single template per
family would leave no variance for that random effect to capture, making
family-level generalization claims unsupported. Model identity is treated
as a stratifying variable (analyses run per model and compared, not
pooled across models as if interchangeable), consistent with RQ2 treating
cross-model heterogeneity as a question in its own right rather than a
nuisance to average away. Male-male and female-female configurations
serve as a within-scenario same-gender control: a genuine agent-gender
effect should appear as an MF/FM asymmetry that is absent (or
substantially smaller) in the MM/FF comparison, which helps rule out
explanations that are really about scenario content rather than gender
per se.

For RQ3, hedge/refusal rate is analyzed as its own binary outcome (hedged
or not, per response) with the same fixed-effect structure as above,
using whatever hedge-coding rule the pilot run settles on
(Section~\ref{sec:protocol}).

A categorical verdict is derived post hoc from \texttt{fault\_rating} via
a cutpoint decided during analysis (Section~\ref{sec:protocol}) and
reported as a secondary, more interpretable view of the same effects, not
as a substitute for the continuous-scale analysis above.

The family-by-family breakdown from RQ2 is additionally read against a
specific, falsifiable prediction from ambivalent-sexism theory
\citep{glickfiske1996,glickfiske1999}, fixed here before any data exist
rather than fit to the results afterward: benevolent sexism toward women
(protective paternalism) predicts female-favoring leniency concentrated
in families framable as the agent needing protection or accommodation
(e.g.\ emotional labor, sexuality and intimacy), while hostility toward
men via resentment of paternalism predicts male-disadvantaging harshness
concentrated in families involving power, control, or a failed
provider/authority role (e.g.\ financial provision, household labor,
jealousy). If the family-by-family results instead show a uniform
direction and magnitude across all 9 families, that argues against an
ambivalent-sexism account and toward a simpler, undifferentiated leniency
or harshness bias -- a distinction the single-family designs in prior work
(Section~\ref{sec:related}) cannot draw.

\subsection{Confirmatory Results}
% TODO -- blocked. No model has been run against the vignette set yet; the
% pilot manipulation/severity check and the confirmatory pass are both
% still pending (owned by a collaborator; tracked separately from this
% drafting pass).
\textcolor{red}{\bf [TODO -- no data collected yet]}
